% chapters/16/section-06.tex
\section{Presentasi Demo, Kesesuaian GDD, dan Refleksi Teknologi}
\label{sec:presentasi-refleksi}

\subsection{Demo Lima sampai Sepuluh Menit}

Demo lima sampai sepuluh menit pada minggu kelima belas, dan uji fungsional
oleh dosen pada minggu keenam belas, merupakan bentuk penilaian yang lazim
pada mata kuliah pengembangan permainan \cite{unima-rps-game}. Waktu sependek
itu menuntut pilihan, dan pilihan yang paling sering salah adalah menceritakan
seluruh proses pembuatan. Yang perlu diperlihatkan bukan perjalanan Anda,
melainkan hasilnya.

Susunan yang dapat dipakai diperlihatkan pada Tabel~\ref{tab:susunan-demo}.

\begin{table}[htbp]
\centering
\caption{Susunan presentasi demo lima sampai sepuluh menit}
\label{tab:susunan-demo}
{\setlength{\tabcolsep}{4pt}\footnotesize%
\begin{tabularx}{\textwidth}{@{}C{1.7cm} L{3.2cm} Y@{}}
\toprule
\headrow \textbf{Waktu} & \textbf{Bagian} & \textbf{Isi} \\
\midrule
1 menit & Pembukaan & Nama permainan, jenisnya, dan satu kalimat tentang apa
yang pemain lakukan di dalamnya. \\
\addlinespace
4--6 menit & Demo langsung & Permainan dijalankan dari awal sampai tujuan, atau
sampai bagian yang paling Anda banggakan. Tidak ada penjelasan bertele-tele di
tengah demo. \\
\addlinespace
1--2 menit & Temuan pengujian & Satu atau dua temuan terpenting dari
\istilah{playtest} antar kelompok, beserta perbaikan yang Anda lakukan. \\
\addlinespace
1 menit & Penutup & Kekurangan yang Anda ketahui, dan satu hal yang akan Anda
kerjakan apabila waktunya bertambah. \\
\bottomrule
\end{tabularx}}
\end{table}

Bagian temuan pengujian adalah bagian yang paling membedakan. Memperlihatkan
bahwa Anda menemukan masalah pada karya Anda sendiri, dan memperbaikinya,
menunjukkan penguasaan yang lebih tinggi daripada memperlihatkan permainan yang
tampak tanpa cacat.

\begin{tipsbox}
Siapkan demo cadangan. Rekam permainan Anda sebagai video pendek, dan bawa
salinan \istilah{build} pada dua tempat. Kegagalan teknis pada saat presentasi
adalah hal yang biasa, dan cara menghadapinya dinilai.
\end{tipsbox}

\subsection{Kesesuaian GDD dengan \istilah{Build}}

GDD Anda ditulis pada awal perkuliahan, ketika sebagian besar hal belum
diketahui. Sekarang \istilah{build} Anda sudah ada, dan perbedaan antara
keduanya hampir pasti ada. Perbedaan itu tidak selalu salah; yang salah adalah
perbedaan yang tidak Anda sadari.

Tiga keadaan yang mungkin Anda temui diperlihatkan pada
Tabel~\ref{tab:gdd-build}.

\begin{table}[htbp]
\centering
\caption{Tiga keadaan kesesuaian antara GDD dan \istilah{build}}
\label{tab:gdd-build}
{\setlength{\tabcolsep}{4pt}\footnotesize%
\begin{tabularx}{\textwidth}{@{}L{3.6cm} Y Y@{}}
\toprule
\headrow \textbf{Keadaan} & \textbf{Artinya} & \textbf{Tindakan} \\
\midrule
Sesuai & Yang dijanjikan GDD benar-benar ada di dalam \istilah{build}. &
Tidak ada yang perlu dikerjakan; cantumkan sebagai bukti pada laporan. \\
\addlinespace
GDD diperbarui & Rancangannya berubah ketika dikerjakan, dan perubahan itu
disengaja. &
Perbarui GDD, dan tuliskan alasan perubahannya. Yang dinilai adalah kesadaran
Anda atas perubahan itu. \\
\addlinespace
Belum dikerjakan & Yang dijanjikan GDD belum ada, dan tidak akan ada. &
Cantumkan sebagai kekurangan yang diketahui. Jangan menghapusnya dari GDD agar
laporannya tampak lengkap. \\
\bottomrule
\end{tabularx}}
\end{table}

Baris ketiga adalah yang paling penting untuk dihindari. Menghapus janji dari
GDD agar hasilnya tampak sesuai adalah bentuk ketidakjujuran akademik, dan
sekaligus menghapus bukti bahwa Anda mampu menilai karya sendiri.

\subsection{Refleksi Pemilihan Teknologi}

Bagian terakhir dari laporan Anda adalah refleksi atas pilihan teknologinya.
Sepanjang mata kuliah ini Anda memakai dua perkakas yang berbeda: Python
beserta Pygame pada Bab~III sampai Bab~XI, dan Unity beserta C\# pada Bab~XII
sampai Bab~XVI. Refleksi itu bukan penilaian mana yang lebih baik, melainkan
keterangan tentang mana yang cocok untuk keperluan apa \cite{rokania-unity}.

Tiga pertanyaan berikut menuntun refleksi Anda, dan jawabannya sebaiknya
disertai contoh dari pekerjaan Anda sendiri, bukan dari buku ini.

\begin{enumerate}
  \item \textbf{Pada bagian mana Pygame mempercepat pekerjaan Anda?} Jawaban
        yang umum adalah pada bagian yang perhitungannya perlu Anda kuasai
        sendiri, dan pada permainan yang aturannya sederhana.
  \item \textbf{Pada bagian mana Unity mempercepat pekerjaan Anda?} Jawaban
        yang umum adalah pada penyusunan level, pengaturan fisika, animasi, dan
        bunyi, sebab pekerjaan itu diurus mesin.
  \item \textbf{Bagian mana yang paling Anda kuasai sekarang, dan bagian mana
        yang masih perlu Anda pelajari?} Jawaban atas pertanyaan ini menjadi
        rencana belajar Anda sesudah mata kuliah ini berakhir.
\end{enumerate}

\begin{catatabox}
Refleksi yang jujur akan memuat hal yang tidak Anda kuasai. Menyebutkannya
bukan pengakuan kelemahan, melainkan bukti bahwa Anda mampu menilai kemampuan
sendiri dengan tepat. Itulah bagian tertinggi pada ranah C5 yang dinilai pada
\SUBCPMK{4.3}: menilai berdasarkan bukti, bukan berdasarkan kesan.
\end{catatabox}

\subsection{Penutup Rangkaian Materi}

Bab ini menutup Materi 14, dan dengannya seluruh Sub-CPMK 4.1 sampai 4.3 pada
\CPMK{4} sudah Anda tuntaskan. Perjalanan Anda sepanjang mata kuliah ini adalah
sebagai berikut.

\begin{table}[htbp]
\centering
\caption{Rangkaian capaian dari Bab III sampai Bab XVI}
\label{tab:rangkaian-capaian}
{\setlength{\tabcolsep}{4pt}\footnotesize%
\begin{tabularx}{\textwidth}{@{}L{3.4cm} L{4.6cm} Y@{}}
\toprule
\headrow \textbf{Bab} & \textbf{Yang Anda bangun} & \textbf{Perkakas} \\
\midrule
Bab~III--XI & Lingkungan, \istilah{game loop}, gerak, fisika, tabrakan, AI
pengejar, mesin keadaan, pencarian jalur, mini proyek &
Python dan Pygame \\
\addlinespace
Bab~XII--XIII & Proyek Unity, level, fisika, kendali pemain &
Unity dan C\# \\
\addlinespace
Bab~XIV--XV & Aset visual, animasi, antarmuka, bunyi, partikel &
Unity dan C\# \\
\addlinespace
Bab~XVI & Pengujian, perbaikan, dan presentasi &
\istilah{Checklist}, \istilah{playtest}, dan pemeriksaan otomatis \\
\bottomrule
\end{tabularx}}
\end{table}

Yang tersisa hanyalah Bab~XVII, yaitu menilai kembali seluruh capaian Anda
terhadap CPL dan CPMK mata kuliah ini, dan menyusun rencana belajar
selanjutnya.
